\section{Uniform growth of the Hermite orbital positivity radius}
\label{sec:analytic-hermite-radius-growth}

The one-variable formula in NN-A030 determines each fixed radius.
A uniform matrix-integral theorem supplies the missing growth rate
when the endpoint multiplicity increases. The matrix-integral
nonvanishing theorem is an imported result; the scaling to the
Hermite radii and to the heat-detection constants is the deduction.

\paragraph{The exact imported theorem and normalization.}
Novak, \emph{Topological Expansion of Oscillatory BGW and HCIZ
Integrals at Strong Coupling}, arXiv:2203.10746v1 (2022), equation
(1.2), p.~2, defines
\[
 I_N(z,A,B)=\int_{U(N)}
          \exp\bigl(zN\operatorname{Tr}(AUBU^*)\bigr)\,dU,
\]
with probability Haar measure and ordinary, unnormalized trace.
The closed polydisc $D_N(\varepsilon)$ is defined on p.~3 by
$|z|\le\varepsilon$ and modulus at most one for each eigenvalue
of $A$ and $B$. The first assertion of Theorem~1.2, p.~6, gives
an absolute $\varepsilon_0\in(0,2/27)$ such that $I_N$ is nonzero
throughout this polydisc for every $N\ge1$. It applies in
particular to the Hermitian matrices used here. The conclusion
is uniform in the dimension, not just a statement for each fixed
dimension. We do not assert a numerical value for $\varepsilon_0$.

\begin{theorem}[NN-A038: square-root growth in multiplicity]
\label{thm:analytic-hermite-radius-growth}
For the Hermite orbital radii $\mathfrak r_m$ of NN-A029 and
$h_m=\max_j|h_{m,j}|$, one has, for every integer $m\ge2$,
\begin{equation}\label{eq:analytic-radius-stable-bounds}
 \frac{2\varepsilon_0 m}{h_m}
       \le\mathfrak r_m\le\frac{\pi(m-1)}{h_m}.
\end{equation}
Consequently
\begin{equation}\label{eq:analytic-radius-root-growth}
 \varepsilon_0\sqrt m\le\mathfrak r_m
                         \le\pi\sqrt{m-1},
 \qquad\mathfrak r_m=\Theta(\sqrt m).
\end{equation}
More precisely,
\begin{equation}\label{eq:analytic-radius-liminf-limsup}
 \varepsilon_0\le
 \liminf_{m\to\infty}\frac{\mathfrak r_m}{\sqrt m}
 \le\limsup_{m\to\infty}\frac{\mathfrak r_m}{\sqrt m}
 \le\frac\pi2.
\end{equation}
These are unconditional finite-dimensional statements. Existence
of a limit in the last display, its value, an explicit admissible
$\varepsilon_0$, and a closed formula for each $m$ are not claimed.
\end{theorem}

\begin{proof}
Let $X=X^*$ have spectral diameter $R>0$, and put
\[
 c=\tfrac12\bigl(\lambda_{\max}(X)+\lambda_{\min}(X)\bigr),
 \quad A=\frac{2(X-cI)}R,\quad B=\frac{D_m}{h_m},
 \quad z=\frac{iRh_m}{2m}.
\]
Both $A$ and $B$ are Hermitian with spectra in $[-1,1]$.
Since $\operatorname{Tr}D_m=0$, invariance of Haar measure and
cyclicity of the trace give exactly
\[
 I_m(z,A,B)
   =\int_{U(m)}e^{i\operatorname{Tr}((X-cI)UD_mU^*)}\,dU
   =\phi_m(X).
\]
If $R\le2\varepsilon_0m/h_m$, the imported theorem makes this
number nonzero. The Hermite spectrum is invariant under negation,
so a unitary permutation conjugates $D_m$ to $-D_m$; hence
$\phi_m(X)$ is real. Along $uX$, $0\le u\le1$, the same
nonvanishing statement applies and the value starts at $1$.
It is therefore strictly positive. A scalar $X$, the case $R=0$,
also has value $1$. The definition of $\mathfrak r_m$ proves the
lower bound in \eqref{eq:analytic-radius-stable-bounds}. The upper
bound is the already proved extended-Chebyshev bound in NN-A030.

For completeness, all required Hermite size estimates follow
from a finite symmetric matrix. Let $J_m$ be the real tridiagonal
$m\times m$ matrix with zero diagonal and off-diagonal entries
$\sqrt1,\sqrt2,\ldots,\sqrt{m-1}$. Expanding its characteristic
determinant along the last row gives
\[
 \det(xI_m-J_m)=x\det(xI_{m-1}-J_{m-1})
                  -(m-1)\det(xI_{m-2}-J_{m-2}),
\]
with initial values $1,x$. This is the defining recurrence of
$\operatorname{He}_m$, so the eigenvalues of $J_m$ are the
$h_{m,j}$. Alternating the signs of the coordinate vectors
conjugates $J_m$ to $-J_m$, and therefore its top eigenvalue is
$h_m=\|J_m\|$. Write $J_m=T_m+T_m^*$, where $T_m$ is the
weighted one-step shift. Then
\[
 h_m\le2\|T_m\|=2\sqrt{m-1},\qquad
 h_m^2\ge\frac1m\operatorname{Tr}(J_m^2)=m-1.
\]
Substitution gives
\[
 \mathfrak r_m\ge\frac{\varepsilon_0m}{\sqrt{m-1}}
                      \ge\varepsilon_0\sqrt m,
 \qquad
 \mathfrak r_m\le\pi\sqrt{m-1},
\]
which proves \eqref{eq:analytic-radius-root-growth}.

We also prove the precise scale of $h_m$, rather than importing
a Hermite asymptotic. For $2\le k<m$, restrict to the last $k$
coordinates. Its off-diagonal weights are at least $\sqrt{m-k}$.
The positive unit eigenvector of the unweighted path on $k$
vertices has entries proportional to $\sin(j\pi/(k+1))$ and
eigenvalue $2\cos(\pi/(k+1))$. Its Rayleigh quotient in this
principal submatrix is thus at least
\[
                  2\sqrt{m-k}\cos\frac\pi{k+1}.
\]
The variational principle gives the same lower bound for $h_m$.
Choose any integer sequence $k\to\infty$ with $k/m\to0$,
and use the preceding upper bound. It follows that
$h_m/(2\sqrt m)\to1$. Dividing
\eqref{eq:analytic-radius-stable-bounds} by $\sqrt m$ proves
\eqref{eq:analytic-radius-liminf-limsup}.
\end{proof}

\paragraph{The resulting heat and representation constants.}
Under A and the exact heat normalization of NN-A029, a positive
endpoint zero $x$ of multiplicity $m\ge2$ contributes
\[
 C_F(x,m)=\frac{x\mathfrak r_m}{2\sqrt2},\qquad
 C_U(x,m)=\frac{x(1-\beta_x)\mathfrak r_m}{2\sqrt2},
 \qquad\beta_x=\frac1{p_0x^2},\quad p_0=F_0'(0).
\]
Thus, uniformly in the abstract multiplicity,
\[
 \frac{\varepsilon_0x\sqrt m}{2\sqrt2}
       \le C_F(x,m)\le\frac{\pi x\sqrt{m-1}}{2\sqrt2},
\]
and the same bounds multiplied by $1-\beta_x>0$ hold for $C_U$.
The least Toeplitz order and index window have the minimum of
the $C_F$ as their exact $\sqrt a$-scaled limit; the least failing
unitary rank has the minimum of the $C_U$, as proved in NN-A029.
These estimates describe the dependence of each candidate on
multiplicity. They do not assert that arithmetic zeros of any
prescribed multiplicity exist or that their multiplicities tend
to infinity.

\paragraph{Geometric and Grothendieck translation.}
The Hermite eigenvalues are distinct, so their conjugacy orbit
is the compact full-flag manifold $U(m)/U(1)^m$, canonically the
complex flag variety of $\mathbb C^m$ as a homogeneous space.
Its specified Haar probability measure has Fourier transform
$\phi_m$; the theorem proves a dimension-uniform positive region
for this concrete transform. The Weyl--HCIZ and Schur dictionary
proved in NN-A028--NN-A030 identifies its spectral-diameter
bound with a bound on the width of every allowed highest-weight
partition in the scaled heat alphabet. Schur and exterior powers
are the actual representation-ring operations of SGA~6,
Expos\'e~V, Example~2.2(b) and Lemma~2.2.1. The displayed
$\Theta(\sqrt m)$ bounds therefore quantify how the local
Artin length $m$ changes these positivity tests. This is a
finite-dimensional flag-orbit statement; no integral polarization
for the nonintegral Hermite weights or arithmetic cohomology
realization is asserted.

\paragraph{Source status and remaining gate.}
The input is the stated theorem of the cited 2022 primary
preprint, which explicitly supersedes arXiv:2006.04304v1.
Its definitions and complete theorem statement were checked,
as was the uniform nonvanishing step in Theorem~4.7,
equations~(4.40)--(4.42), p.~34. That step first treats all large
$N$ and then shrinks the disc to handle the finitely many smaller
dimensions. We have not claimed to reproduce or independently
verify all of its combinatorial proof. A sharper remaining
question is whether $\mathfrak r_m/\sqrt m$ converges, and, if
so, its value and a quantitative error. The growth exponent is
resolved by the displayed deduction.
